% Introducción (ejemplo): presenta el contexto, la motivación y los objetivos.
En los artículos originales enviados a la revista IEEE, la introducción debe presentar de forma concisa el contexto y la relevancia del problema abordado, describiendo brevemente el área de estudio y su importancia en el campo de la ingeniería o tecnología; incluir un resumen del estado del arte con referencia a los trabajos más relevantes y recientes, destacando las limitaciones o vacíos existentes que motivan la investigación; exponer con claridad el objetivo y la principal contribución del trabajo, indicando qué se propone y en qué se diferencia o mejora respecto a lo ya publicado; y, opcionalmente, cerrar con una breve descripción de la estructura del artículo para guiar al lector en el desarrollo del contenido.

% Mantén párrafos cortos y evita resultados aquí.
\textcolor{blue}{La escritura académica en LaTeX ofrece control tipográfico y reproducibilidad. Esta plantilla en español muestra, con ejemplos mínimos, cómo estructurar un artículo usando IEEEtran: figuras, tablas, ecuaciones y bibliografía. Úsala para practicar el flujo básico de edición y compilación.}


% Estructura sugerida del artículo
El resto del documento se organiza así: en la Sección~\ref{sec:experimental} se describe un ejemplo de desarrollo experimental simplificado; en la Sección~\ref{sec:metodologia} se explica la metodología con ecuaciones de ejemplo; en la Sección~\ref{sec:resultados} se muestran tablas y figuras de ejemplo; y en la Sección~\ref{sec:conclusiones} se listan conclusiones y trabajo futuro.


% Ejemplo de cita para mantener algunas referencias en el .bib
\textcolor{red}{LaTeX facilita gestionar referencias con BibTeX, por ejemplo \cite{Cooley1965,Geron2019}. También puedes usar estilos numéricos compatibles con IEEE. Las citas deben registrarse en el archivo reference.bib}